\section{Preliminaries}\label{sec:prelim}

Throughout, $\A=\{0,1\}^2$, and a cell $(x,y)\in\A$ is also written as the integer $2x+y\in\{0,1,2,3\}$; thus $3=(1,1)$ and $0=(0,0)$. The left shift is $\sh(c)_i=c_{i+1}$. A \emph{block map of radius $R$} is a map $g\colon\A^{\Z}\to\A^{\Z}$ of the form $g(c)_i=\gamma(c_{i-R},\dots,c_{i+R})$ for a local rule $\gamma\colon\A^{2R+1}\to\A$; we identify $g$ with the lookup table of $\gamma$, which has $4^{2R+1}$ entries of two bits each. A block map of radius $R$ is also a block map of every larger radius. By Curtis--Hedlund--Lyndon, $\End(\A^{\Z})$ is the monoid of block maps and $\Aut(\A^{\Z})$ the group of invertible ones. We write
\[
\Cent_R(F)=\{g\ \text{a block map of radius}\le R:\ g\circ F=F\circ g\},\qquad \Cent(F)=\bigcup_R\Cent_R(F),
\]
for the centralizer of $F$ in $\End(\A^{\Z})$, and $\Cent^{\times}_R(F)$, $\Cent^{\times}(F)$ for the invertible elements. For a configuration $c$ and integers $i<j$ we write $c_{[i,j)}=c_ic_{i+1}\cdots c_{j-1}$, and $c_{[i,j]}=c_{[i,j+1)}$.

For an elementary rule $f$, identified with its Wolfram number $\sum_{a,b,c}f(a,b,c)2^{4a+2b+c}$, the automaton $F_f$ is defined by \eqref{eq:F}. It has radius $1$, and $F_f^b$ has radius at most $|b|$, so $\sh^aF_f^b\in\Cent_{|a|+|b|}(F_f)$.

\begin{definition}
The \emph{spacetime subshift} of a reversible cellular automaton $F$ on $\A^{\Z}$ is
\[
X_F=\bigl\{(F^n(c))_{n\in\Z}:\ c\in\A^{\Z}\bigr\}\subset\A^{\Z^2},
\]
with the $\Z^2$-action generated by the horizontal shift $\sh$ and the vertical shift $\tau$, $(\tau z)_{n}=z_{n+1}$. An \emph{automorphism} of $X_F$ is a homeomorphism of $X_F$ commuting with both.
\end{definition}

\begin{lemma}\label{lem:spacetime}
For reversible $F$ the map $c\mapsto(F^n(c))_{n\in\Z}$ is a homeomorphism $\A^{\Z}\to X_F$ conjugating $(\sh,F)$ to $(\sh,\tau)$, and it induces an isomorphism
\[
\Aut(X_F)\;\cong\;\Cent^{\times}(F)=\{g\in\Aut(\A^{\Z}):\ gF=Fg\}.
\]
Under this isomorphism $\langle\sh,\tau\rangle$ corresponds to $\langle\sh,F\rangle$.
\end{lemma}

\begin{proof}
The map is continuous, injective (the $0$-th row recovers $c$), and surjective onto $X_F$ by definition, hence a homeomorphism of compact spaces; it carries $\sh$ to $\sh$ and $F$ to $\tau$ since $F^{n}(F(c))=F^{n+1}(c)$. A homeomorphism of $\A^{\Z}$ commuting with $\sh$ is a block map by Curtis--Hedlund--Lyndon, and commuting with $\tau$ means commuting with $F$.
\end{proof}

So determining $\Aut(X_F)$ means determining which invertible block maps commute with $F$. The non-invertible elements of $\Cent(F)$ are the endomorphisms of $X_F$ that are not automorphisms; the constant map onto a fixed point of $F$ is always one. For every rule with $f(0,0,0)=f(1,1,1)=0$, which includes $54$ and $90$, both $(0,0)^{\Z}$ and $(1,1)^{\Z}$ are fixed by $F_f$; for rule $150$, which has $f(1,1,1)=1$, only $(0,0)^{\Z}$ is.

\begin{sloppypar}
\emph{Conjugacies.} Let $\mathrm{rev}(c)_i=c_{-i}$ be the reflection and $\kappa(x,y)=(1-x,1-y)$ the cellwise complement. By \eqref{eq:F}, $\mathrm{rev}\,F_f\,\mathrm{rev}=F_{f'}$ with $f'(a,b,c)=f(c,b,a)$, and $\kappa F_f\kappa=F_{f''}$ with $f''(a,b,c)=f(1-a,1-b,1-c)$. Both $\mathrm{rev}$ and $\kappa$ are involutions that act cellwise up to reflection, so conjugation by them maps $\Cent_R(F_f)$ bijectively onto $\Cent_R(F_{f'})$, respectively $\Cent_R(F_{f''})$, preserves invertibility and maps $\langle\sh,F_f\rangle$ onto the corresponding group. So $\Aut(X_{F_f})$, and every statement below about block maps of bounded radius, depends only on the class of $f$ under these two operations. There are $88$ classes \cite{Takesue1987}, \cite[Table~I]{Takesue1989}; one representative of each appears in Table~\ref{tab:walls}. Wolfram's complement $1\oplus f(1-a,1-b,1-c)$ is not a conjugacy of second-order automata, so these are not the classes of the elementary rules \cite[p.~377 and footnote~3]{Takesue1989}.
\end{sloppypar}

\subsection{The group generated by the shift and the automaton}

The automorphisms constructed below lie outside $\langle\sh,F_f\rangle$. To make statements about the group they generate together with $\sh$ and $F_f$ we need the structure of $\langle\sh,F_f\rangle$ itself. For $c\in\A^{\Z}$ let $Y_c(t,i)$ be the $y$-bit of $F^t(c)_i$; then the $x$-bit of $F^t(c)_i$ is $Y_c(t-1,i)$, and by \eqref{eq:F}
\[
Y(t+1,i)\oplus Y(t-1,i)=f\bigl(Y(t,i-1),Y(t,i),Y(t,i+1)\bigr)\qquad\text{for all }(t,i).
\]
A block map $g$ \emph{depends on} a bit of the cell $c_i$ (for the output bit $y(g(c)_0)$, say) if two configurations that differ only in that bit have different outputs. For a nonconstant block map $g$ let $r^+(g)$ be the largest and $-r^-(g)$ the smallest $i$ such that $g(c)_0$ depends on $c_i$. Then $r^\pm(\sh^ag)=r^\pm(g)\pm a$, and $r^\pm(F_f^n)\le|n|$ because $F_f^n(c)_0$ depends only on $c_{[-|n|,|n|]}$.

\begin{lemma}[free staircases]\label{cl:lem:stair}
Let $f$ be any elementary rule and $n\ge1$. For every assignment of bits to the $2n$ points $(k,n-k-2)$ and $(k,n-k-1)$, $0\le k<n$, there is a configuration $c$ whose diagram $Y_c$ takes the assigned values at these points.
\end{lemma}

\begin{proof}
The points lie on the two diagonals $d_m=\{(t,i):t+i=m\}$ with $m=n-2$ and $m=n-1$. Put $K=2n+2$, $\tau(m)=m-n+1+K$, and let $\Omega$ be the set of $(t,i)$ with $n-1-K\le t+i\le n-1$ and $-1\le t\le\tau(t+i)$; so $\Omega\cap d_m=\{(t,m-t):-1\le t\le\tau(m)\}$ for $n-1-K\le m\le n-1$, and $\tau(m+1)=\tau(m)+1$. We define $Y\colon\Omega\to\{0,1\}$, one diagonal at a time, from right to left. On $d_{n-1}$ and $d_{n-2}$ let $Y$ take the assigned values at the $2n$ points, which lie in $\Omega$ because $k\le n-1<\tau(n-2)$, and the value $0$ elsewhere. For $m=n-3,n-4,\dots,n-1-K$ in turn put $Y(\tau(m),m-\tau(m))=0$ and then, for $t=\tau(m),\tau(m)-1,\dots,0$,
\begin{equation}\label{cl:eq:fill}
Y(t-1,m-t+1)=Y(t+1,m-t+1)\oplus f\bigl(Y(t,m-t),\,Y(t,m-t+1),\,Y(t,m-t+2)\bigr).
\end{equation}
The point $(t,m-t)$ of $d_m$ has just been defined, and $(t,m-t+1)\in d_{m+1}$, $(t+1,m-t+1)$ and $(t,m-t+2)\in d_{m+2}$ have been defined before, because $t\le\tau(m)<\tau(m+1)$ and $t+1\le\tau(m+2)$. So every point of $\Omega$ receives a value, and \eqref{cl:eq:fill} is the spacetime relation at the point $(t,m-t+1)$.

Conversely, suppose that the five points $(t\pm1,i)$, $(t,i\pm1)$ of the relation at $(t,i)$ all lie in $\Omega$, and put $m=t+i-1$. Then $(t,i-1)\in d_m$ and $(t,i+1)\in d_{m+2}$ give $n-1-K\le m\le n-3$ and $t\le\tau(m)$, and $(t-1,i)\in\Omega$ gives $t\ge0$. So the relation at $(t,i)$ is \eqref{cl:eq:fill} for this $m$ and $t$, and it holds. Thus $Y$ satisfies every instance of the spacetime relation whose five points lie in $\Omega$.

Let $c$ be any configuration with $c_i=(Y(-1,i),Y(0,i))$ for $n-K\le i\le n-1$; these points lie in $\Omega$. We claim that $Y_c(t,i)=Y(t,i)$ for $t\in\{-1,0\}$ and $n-K\le i\le n-1$, and for $t\ge1$ and $n-K+t\le i\le n-1-t$. For $t\le0$ this is the choice of $c$. For $t\ge1$, $Y_c(t,i)=Y_c(t-2,i)\oplus f(Y_c(t-1,i-1),Y_c(t-1,i),Y_c(t-1,i+1))$, the four points on the right lie in the ranges already covered for the rows $t-2$ and $t-1$, and the five points of the relation at $(t-1,i)$ lie in $\Omega$ (their diagonals are between $t+i-2\ge n-K$ and $t+i\le n-1$, and their rows between $t-2\ge-1$ and $t\le\tau(t+i)$, the last since $i\ge n-K+t>n-1-K$); so $Y$ satisfies the same relation there, and the claim follows by induction on $t$. The assigned points $(k,n-k-2)$, $(k,n-k-1)$ lie in these ranges: for $k=0$ in row $0$, and for $k\ge1$ because $n-K+k\le n-k-2$, that is $K\ge2k+2$. So $Y_c$ takes the assigned values there.
\end{proof}

\begin{theorem}\label{cl:thm:A}
Let $f$ be an elementary rule, $F=F_f$ and $n\ge1$.
\begin{enumerate}
\item If $f$ depends on its third argument, then $y(F^n(c)_0)$ depends on $y(c_n)$, and $x(F^{-n}(c)_0)$ depends on $x(c_n)$. So $r^+(F^n)=r^+(F^{-n})=n$.
\item If $f$ depends on its first argument, then $y(F^n(c)_0)$ depends on $y(c_{-n})$, and $x(F^{-n}(c)_0)$ depends on $x(c_{-n})$. So $r^-(F^n)=r^-(F^{-n})=n$.
\item The map $(a,b)\mapsto\sh^aF^b$ is an isomorphism $\Z^2\to\langle\sh,F\rangle$ if and only if $f$ is not a function of its middle argument alone.
\end{enumerate}
\end{theorem}

\begin{proof}
(1) Choose $(a^*,b^*)$ with $f(a^*,b^*,0)\ne f(a^*,b^*,1)$. By Lemma~\ref{cl:lem:stair} there is a configuration $c$ with $Y_c(k,n-k-2)=a^*$ and $Y_c(k,n-k-1)=b^*$ for $0\le k<n$. Let $c'$ be $c$ with $y(c_n)$ complemented, and $\Delta=Y_c\oplus Y_{c'}$. We show by induction on $k$, $0\le k\le n$, that $\Delta(k,i)=0$ for $i<n-k$, $\Delta(k,n-k)=1$, and $\Delta(k-1,i)=0$ for $i<n-k+1$. For $k=0$ this holds because $c$ and $c'$ differ only in $y(c_n)$. For the step from $k<n$ to $k+1$, the relation at $(k,i)$ gives $\Delta(k+1,i)=0$ for $i<n-k-1$, since all four inputs $\Delta(k-1,i)$, $\Delta(k,i-1)$, $\Delta(k,i)$, $\Delta(k,i+1)$ vanish. At $i=n-k-1$ the inputs $\Delta(k-1,n-k-1)$, $\Delta(k,n-k-2)$ and $\Delta(k,n-k-1)$ vanish, so both diagrams have the values $a^*,b^*$ at $(k,n-k-2)$, $(k,n-k-1)$, while $\Delta(k,n-k)=1$; hence $\Delta(k+1,n-k-1)=f(a^*,b^*,0)\oplus f(a^*,b^*,1)=1$. For $k=n$ this says $Y_c(n,0)\ne Y_{c'}(n,0)$, that is $y(F^n(c)_0)\ne y(F^n(c')_0)$. For $F^{-n}$ use the time reversal $\rho(x,y)=(y,x)$, which satisfies $\rho F\rho=F^{-1}$ for every $f$ \cite[p.~2500]{Takesue1987}: $x(F^{-n}(c)_0)=y(F^n(\rho c)_0)$ depends on $y((\rho c)_n)=x(c_n)$. With $r^+(F^{\pm n})\le n$ this gives $r^+(F^{\pm n})=n$.

(2) The reflection conjugates $F_f$ to $F_{f'}$ with $f'(a,b,c)=f(c,b,a)$, and $f'$ depends on its third argument; apply (1) to $f'$.

(3) Suppose $f$ depends on its third argument. By (1), $r^+(\sh^aF^b)=a+|b|$ for all $(a,b)$. If $\sh^aF^b=\mathrm{id}$, then also $\sh^{-a}F^{-b}=\mathrm{id}$, so $a+|b|=0=-a+|b|$ and $a=b=0$. Since $\sh$ and $F$ commute, $(a,b)\mapsto\sh^aF^b$ is a homomorphism onto $\langle\sh,F\rangle$, and it is injective. If $f$ depends on its first argument, the same argument with $r^-(\sh^aF^b)=|b|-a$ applies. If $f$ depends on neither, then $f(a,b,c)=g(b)$, and $F(c)_i=(y_i,x_i\oplus g(y_i))$ applies one permutation of $\A$ to every cell, so $F$ has finite order and the map is not injective.
\end{proof}

\begin{proposition}\label{cl:prop:kappa}
For $f\in\{126,219\}$ the cellwise complement $\kappa$ commutes with $F_f$ and $\kappa\notin\langle\sh,F_f\rangle$, so $\langle\sh,F_f,\kappa\rangle=\langle\sh,F_f\rangle\times\langle\kappa\rangle\cong\Z^2\times\Z/2$.
\end{proposition}

\begin{proof}
Since $\kappa F_f\kappa=F_{f''}$ with $f''(a,b,c)=f(1-a,1-b,1-c)$, $\kappa$ commutes with $F_f$ if and only if $f(1-a,1-b,1-c)=f(a,b,c)$ for all $a,b,c$, that is, if and only if $f$ is self-conjugate in the sense of \cite[Appendix~A]{Takesue1989}, which holds for $126$ (which vanishes exactly on $000$ and $111$) and for $219$ (which vanishes exactly on $010$ and $101$). For $126$, $f(0,0,0)=0$, so $(0,0)^{\Z}$ is fixed by $\sh$ and $F_f$, hence by every element of $\langle\sh,F_f\rangle$, while $\kappa$ maps it to $(1,1)^{\Z}$. For $219$, the orbit of $(01)^{\infty}$ under $\langle\sh,F_{219}\rangle$ consists of the four configurations of period $2$ with words $01,10,02,20$, and $\kappa$ maps $(01)^{\infty}$ to $(32)^{\infty}$, which is not among them. Since an element of $\langle\sh,F_f\rangle$ maps each configuration into its orbit, $\kappa\notin\langle\sh,F_f\rangle$. Both rules depend on their first and third arguments, so $\langle\sh,F_f\rangle\cong\Z^2$ by Theorem~\ref{cl:thm:A}(3). As $\kappa$ is an involution commuting with $\sh$ and $F_f$, and $\langle\sh,F_f\rangle\cong\Z^2$ is torsion-free, $\langle\kappa\rangle\cap\langle\sh,F_f\rangle=\{\mathrm{id}\}$, which gives the direct product.
\end{proof}
